% Physical unit 041; accepted Session 03 after fresh 400-dpi authority replay.
\GordanUnitStart{041}{printed page 36}
\GordanTopFolio{36}
\begingroup\small
\[
=\arc\left(\cos={\cos l_0\over\sqrt{{\sin^2l_0\over\sin^2l_2}-\sin^2l_0}}\right)-\arc\left(\cos={\cos l_0\over\sqrt{{\sin^2l_0\over\sin^2l_1}-\sin^2l_0}}\right)
\]
\[
=\arc(\cos=\cot l_0\tg l_2)-\arc(\cos=\cot l_0\tg l_1)
\]
quam aequationem ad elegantiorem formam reducere possumus usi notissima illa identitate:
\[
\left|\begin{array}{ccc}
1&\cos x&\cos y\\
\cos x&1&\cos(x-y)\\
\cos y&\cos(x-y)&1
\end{array}\right|
\]
in qua posito:
\[
 x=\cot l_0\tg l_2,
 \qquad y=\cot l_0\tg l_1,
 \qquad x-y=\psi+\chi
\]
fit:
\[
\left|\begin{array}{ccc}
1&\tg l_2\cot l_0&\tg l_1\cot l_0\\
\tg l_2\cot l_0&1&\cos(\psi+\chi)\\
\tg l_1\cot l_0&\cos(\psi+\chi)&1
\end{array}\right|
=
\left|\begin{array}{ccc}
\tg^2l_0&\tg l_2&\tg l_1\\
\tg l_2&1&\cos(\psi+\chi)\\
\tg l_1&\cos(\psi+\chi)&1
\end{array}\right|=0
\]
vel:
\[
\text{I.}\qquad \tg^2l_0\sin^2(\psi+\chi)=\tg^2l_1+\tg^2l_2-2\tg l_1\tg l_2\cos(\psi+\chi).
\]
Qua in aequatione $l_0$ occurrit et in \gq{$\tg^2l_0$} et in \gq{$\chi$}. Sin hoc temporis momento non stricte libet definire $l_0$, sed valorem quendam $l_0'$, qui ipsi propinquus est; $\chi$ in formula negligi potest, quod continens factorem $e^2$ minimum est.

Ergo $l_0'$, definitum e formula:
\[
\text{II.}\qquad \tg^2l_0'\sin^2\psi=\tg^2l_1+\tg^2l_2-2\tg l_1\tg l_2\cos\psi
\]
differt a $l_0$ quantitate minima secundi gradus, si comparas cum excentricitate.

Ad computandum $l_0''$, propius accedens valore suo ad $l_0$, non plane abjiciendum est $\chi$, sed habenda ratio est membri
\[
 {e^2\over2}\cos l_0(\varphi_2-\varphi_1)
\]
seriei, secundum dignitates quadrati $e^2$ progredientis, in quam integrale $\chi$ explicari potest.
Tum valor $l_0''$, ex aequatione hac definitus:
\[
\text{III.}\quad \tg^2l_0''\sin^2\left[\psi+{e^2\over2}\cos l_0''(\varphi_2-\varphi_1)\right]
=\tg^2l_1+\tg^2l_2-2\tg l_1\tg l_2\cos\left[\psi+{e^2\over2}\cos l_0''(\varphi_2-\varphi_1)\right]
\]
a $l_0$ recedit quantitate, quae, comparata cum excentricitate, minima quarti gradus est.

Primo obtutu putes graves difficultates obstare computationi quantitatis $l_0''$ ex aequatione III. Si quidem $l_0''$ non modo in $\tg^2l_0$ sed etiam in:
\[
 {e^2\over2}\cos l_0''(\varphi_2-\varphi_1)
\]
\endgroup
